% Sendable version of doc/discrepancies.md, typeset in the paper's notation.
% Build:  scripts/make-discrepancies-pdf.sh   (pdflatex, twice)
% KEEP IN SYNC with doc/discrepancies.md -- that is the repository-facing copy,
% this one is what gets sent to the authors.
\documentclass[11pt,a4paper]{article}

\usepackage[T1]{fontenc}
\usepackage[utf8]{inputenc}
\usepackage{lmodern}
\usepackage{textcomp}
\usepackage[margin=2.6cm]{geometry}
\usepackage{amsmath,amssymb}
\usepackage{booktabs}
\usepackage{array}
\usepackage{enumitem}
\usepackage[colorlinks=true,linkcolor=black,urlcolor=blue]{hyperref}

\renewcommand{\arraystretch}{1.2}

% The paper's notation.
\newcommand{\RN}{\mathrm{RN}}
\newcommand{\ulp}{\mathrm{ulp}}
\newcommand{\uls}{\mathrm{uls}}
\newcommand{\ufp}{\mathrm{ufp}}
\newcommand{\alg}[1]{\mathit{#1}}
\newcommand{\Prodacc}[1]{\alg{3Prod}^{\mathrm{acc}}_{#1}}
\newcommand{\Prodfast}[1]{\alg{3Prod}^{\mathrm{fast}}_{#1}}

\title{Formal verification of \emph{Algorithms for triple-word arithmetic}\\[2pt]
       \large verified results, and where they differ from the paper}
\author{Laurent Théry\\\small Inria \quad\texttt{laurent.thery@inria.fr}}
\date{29 July 2026}

\begin{document}
\maketitle

\noindent
This note summarises a complete machine-checked verification of the paper, and
lists every point where the formal proof departs from the printed text. It uses
the paper's own notation throughout; no knowledge of the proof assistant is
needed to read it.

Everything below has been checked by machine. Counterexamples are exact and
reproducible; where a constant differs, both values are proved.

\section{What has been verified}

Algorithms 1--15, together with the appendix Algorithms 18 and 20, and
Theorems 1--11. Radix~2, precision~$p$, round to nearest ties-to-even,
$u = 2^{-p}$.

\begin{center}
\begin{tabular}{@{}lll@{}}
\toprule
 & published & verified \\
\midrule
Thm~1, Cor.~1 ($\alg{VecSum}$)   & ---                       & as stated \\
Thm~2 ($\alg{VSEB}$)             & ---                       & as stated \\
Thm~3 (keep first $k$)           & $2u^k + 4.2u^{k+1}$       & as stated, for general $k$ \\
Thm~4 ($\alg{ToTW}$)             & $p \geq 4$                & as stated \\
Thm~5 ($\alg{RoundTW}$)          & $\alg{RoundTW}(\bar x) = \RN(\bar x)$
                                 & as stated, \textbf{after correcting Alg.~7} (\S\ref{sec:roundtw}) \\
Thm~6 ($\alg{TWSum}$)            & $p \geq 4$
                                 & as stated (\textbf{its proof has a false step}, \S\ref{sec:thm6}) \\
Thm~7 (Alg.~9)                   & $28u^3 + 107u^4$          & as stated \\
Thm~8 (Alg.~11)                  & $10.5u^3 + 39u^4$
                                 & as stated (\textbf{gap in the case analysis}, \S\ref{sec:thm8}) \\
Thm~9 (Alg.~13)                  & $11.5u^3 + 1465u^4$       & $11.5u^3 + \mathbf{1830}u^4$ \\
                                 & $19u^3 + 1502u^4$         & $19u^3 + \mathbf{1870}u^4$ \quad (\S\ref{sec:delta2}) \\
Thm~10 (Alg.~14)                 & $24u^3 + 1509u^4$         & $\mathbf{27}u^3 + \mathbf{2500}u^4$ \\
                                 & $39u^3 + 1582u^4$         & $\mathbf{42}u^3 + \mathbf{2575}u^4$ \quad (\S\ref{sec:delta3}) \\
Thm~11 (Alg.~15)                 & $24u^3 + 10260u^4$        & \textbf{as stated} \\
                                 & $39u^3 + 10333u^4$        & \textbf{as stated} \quad (\S\ref{sec:thm11}) \\
\bottomrule
\end{tabular}
\end{center}

\medskip
\noindent
Two of the $u^3$ constants move (Theorem~10) and two of the $u^4$ constants
move (Theorem~9). \textbf{Every other $u^3$ constant in the paper is confirmed
exactly} --- including $11.5 = 10.5 + 1$ and $19 = 18 + 1$, and all of
Theorem~11.

Not formalised: Algorithms 16, 17, 19 (the sign-folded variants --- they save
one operation each and change no bound).

\section{Points that need a correction}

\subsection{Algorithm 7 ($\alg{RoundTW}$): the inner test has the wrong
polarity}\label{sec:roundtw}

As printed, the first \textbf{if} returns $\RN(x_0+x_1)$ when
\[
  \RN(x_0 + 2x_1) \text{ is inexact}
  \quad\textbf{or}\quad
  \RN\!\left(-\left(\tfrac32 u - 2u^2\right)\!\cdot x_0\right) \neq x_1 .
\]
The second disjunct should be ``$=$'', not ``$\neq$''.

At a genuine midpoint $x_1 = \pm\frac12\ulp(x_0)$, the quantity
$\RN(-(\frac32 u - 2u^2)x_0)$ has magnitude $\approx \frac32 u|x_0|$ and, for
positive $x_1$, the opposite sign --- so it never equals $x_1$, the printed
test is true, and the algorithm returns $\RN(x_0+x_1)$, \textbf{ignoring $x_2$
exactly when $x_2$ decides the rounding}. Conversely, the paper's own
``special case'' $x_0 = 1+2u$, $x_1 = -\frac32 u$ is precisely where
$\RN(-(\frac32 u - 2u^2)x_0) = x_1$, so the printed condition takes the
directional branch there --- also the wrong way round.

\paragraph{Counterexample (binary64).}
$\bar x = (1, u, u^2)$ is a valid triple word, and $x_0 + x_1$ is exactly the
midpoint of $[1,\, 1+2^{-52}]$.

\begin{center}
\begin{tabular}{@{}ll@{}}
\toprule
$\RN(\bar x)$                    & $1 + 2^{-52}$ \\
Algorithm 7 as printed           & $1$ \quad --- wrong \\
Algorithm 7 with ``$=$''         & $1 + 2^{-52}$ \\
\bottomrule
\end{tabular}
\end{center}

\noindent
An exact sweep over random valid triple words fails on about $0.3\%$ of
inputs (all of them midpoints) for $p = 4, 6, 8, 10$. Theorem~5 is true of the
corrected algorithm, and the proof text then reads correctly as written.

\subsection{Lemma 1: two missing hypotheses}

Lemma~1 ($\frac12\ulp(x)$ divides $\RN(x+y)$, used in the proof of Theorem~7)
needs $y$ to be a floating-point number and $x \neq 0$. It is false without
them.

\subsection{Theorem 6: an intermediate claim of the proof is
false}\label{sec:thm6}

The proof asserts that the output of $\alg{VecSum}$ is F-nonoverlapping. It is
not:
\[
  (15,\; 15,\; \tfrac{15}{16},\; \tfrac{15}{16}) \quad\text{at } p = 4
\]
is a legal input whose $\alg{VecSum}$ output violates Definition~2.
\textbf{Theorem~6 itself is true} --- $\alg{VSEB}$ repairs the defect --- but
the argument has to go through $\alg{VSEB}$ rather than through
$\alg{VecSum}$ alone.

\subsection{Theorem 8: a gap in the case analysis}\label{sec:thm8}

The case analysis splits on $\bar x\bar y \geq \frac32 - 5u$ and
$\bar x\bar y < \frac32 - 6u$, leaving the interval between the two uncovered.
Using the single threshold $\frac32 - 7u$ closes it, and the proof goes through
unchanged. \textbf{The statement is unaffected.}

\subsection{Section 8.3: the tail bound needs ulp, not uls}

The argument that the leading limb of $\Prodacc{2,3}(\bar b, \bar x)$ equals
$1$ bounds the tail using F-nonoverlapping, that is, using $\uls$. That is not
enough:
\[
  (16,\; -8,\; -4,\; -2,\; -1)
\]
is F-nonoverlapping, sums to $1$, and has leading limb $16$. What the argument
actually needs is the \textbf{$\ulp$-scale} separation of the two leading
limbs --- which the $\alg{2Sum}$ structure of the inner $\alg{VecSum}$ does
supply, so the conclusion $e_0 = 1$ stands. Only the justification has to
change.

\subsection{$\delta_2$ (Algorithm 18) is $u^3 + 620u^4$, not
$u^3 + 260u^4$}\label{sec:delta2}

Two distinct issues.

\paragraph{(a) Remark 10 is false.}
Algorithm~18 uses $\alg{Fast2Sum}(b_1, z_{01}^{+})$ without establishing the
ordering, on the grounds that when the condition fails $b_1$ is negligible.
But the condition fails exactly when $|b_1| < |z_{01}^{+}|$, and the worst case
$|b_1| \approx \ulp(z_{01}^{+})$ loses $\approx u\,|z_{01}^{+}| \approx 41u^3$
--- not $O(u^4)$. Legal binary64 inputs reach $\mathbf{32u^3}$, that is,
$32$ times the announced $\delta_2$. The last line,
$\alg{Fast2Sum}(e_1, e_2)$, has the same hole ($|e_1| < |e_2|$ is legal, and
costs $\approx 3u^3$).

\emph{This is repairable at no arithmetic cost}: sort the two arguments first.
A $\alg{Fast2Sum}$ preceded by one test is error-free unconditionally, so
Algorithm~18 keeps its 20 operations and gains 2 tests.

\paragraph{(b) The $u^4$ term.}
With the repair in place we prove $\delta_2 = u^3 + 620u^4$. The $u^3$ term is
exactly the announced one; the $u^4$ term is not. Consequence for Theorem~9:

\begin{center}
\begin{tabular}{@{}lll@{}}
\toprule
                & published            & verified \\
\midrule
accurate        & $11.5u^3 + 1465u^4$  & $11.5u^3 + 1830u^4$ \\
fast            & $19u^3 + 1502u^4$    & $19u^3 + 1870u^4$ \\
\bottomrule
\end{tabular}
\end{center}

\subsection{The exponent range}

All the results are proved with an unbounded exponent range --- the model the
paper works in implicitly. This is worth stating in the text, because
Theorem~1 and Corollary~1 are genuinely false near the underflow threshold.

\section{One place where our bound is weaker, and we believe the paper is
right}\label{sec:delta3}

\subsection{$\delta_3$ (Algorithm 20) and Theorem 10}

We can only prove $\delta_3 = 6u^3 + 1250u^4$ against the announced
$3u^3 + 264u^4$, whence Theorem~10 at $27u^3 + 2500u^4$ and $42u^3 + 2575u^4$
instead of $24u^3 + 1509u^4$ and $39u^3 + 1582u^4$.

The reason is that a \emph{triple} word only gives $|x_1| < \ulp(x_0)$ and
$|x_2| < \ulp(x_1)$, that is, twice a double word's separation, so the two
neglected roundings cost $2u^3$ and $4u^3$ instead of $1u^3$ and $2u^3$. The
announced $3u^3$ is what the double-word separations would give.

\textbf{We do not claim this is an error.} A numerical search over legal inputs
reaches only about $2.5u^3$, so $3u^3$ does look attained and the remaining
factor of two is very likely slack in our proof, not in yours. We record it
only so that the discrepancy is not mistaken for a disagreement.

\subsection{A hypothesis that is implicit in Algorithm 20}

Worth making explicit in the appendix: $\delta_3$ depends on \textbf{how close
the second factor is to $1$}, not merely on its leading limb being $1$.
Algorithm~14 supplies a second factor within $40u^2$ of $1$; Algorithm~15 can
only supply one within $\approx 101u^2$, because there the seed error enters
undamped (see \S\ref{sec:thm11}). Our bound is therefore stated as a function
of that tolerance $c$:
\[
  \delta_3(c) = 6u^3 + (31c + 10)\,u^4
  \qquad\text{for a second factor within } c\,u^2 \text{ of } 1 .
\]
The $u^3$ term does not depend on $c$ at all.

\section{Theorem 11 (Algorithm 15)}\label{sec:thm11}

\subsection{It holds exactly as published}

$24u^3 + 10260u^4$ and $39u^3 + 10333u^4$ at $p \geq 11$, verified --- and this
is the only theorem in the paper for which the printed constants survive
despite \emph{every} intermediate constant of ours being worse ($\delta_3$
above; and we could only reach $100u^2$ for the seed, where the supplementary
material proves $81u^2 + 622u^3$).

That is not luck. \textbf{Our proof takes a different route}, and has to. The
supplementary material bounds $|\delta_1 - (\delta_1+\delta_2)/2|$ by the
triangle inequality, which gives $\delta_1$ the weight $1.5$. That discards a
real cancellation: $i^{(1)}$ occurs twice --- once as a factor of the result,
once inside $i^{(2)}$ --- with opposite signs. Keeping it gives $\delta_1$ the
weight $1/2$, and the total becomes
\[
  \left|y - \sqrt{\bar x}\right| \;\leq\;
  \left(\tfrac12\delta_1 + \tfrac12\delta_2 + \delta_3 + 15200u^4\right)
  \sqrt{\bar x},
\]
which is what leaves enough $u^3$ room to absorb our worse $u^4$ terms. Done
termwise, with our $\delta_3$, Theorem~11 would come out at $29u^3$ / $44u^3$
and the published bound would fail on both terms at once.

This suggests that the weight $1.5$ on $\delta_1$ is simply loose, and that the
true first-order sensitivity of Algorithm~15 to its first product is one half,
not one and a half.

\subsection{$p \geq 11$ is confirmed, and confirmed sharp}

The supplementary material's $9916u^4$ residual is exactly $E^2(3+E)/2$
evaluated at $E = 81u^2 + 622u^3$: that is $9915.5u^4$ at $p = 11$, and
$9989.9u^4$ at $p = 10$, where it exceeds the stated value. So the precision
requirement is real, the constant is what it is announced to be, and the two
are consistent.

\section{The reference implementation}\label{sec:impl}

We also read the C reference implementation of this arithmetic
(\texttt{triple\_float.c} from \texttt{NXP-Research/TWFalcon}, binary32). It
matters here because it is independent evidence about the two points above.

\paragraph{The algorithms are transcribed faithfully.}
\texttt{tw\_sqrt} follows Algorithm~15 line for line --- the $1+4u$ seed, the
exact halvings, $\mathtt{1.5f} - h_{01}^{(2)}$, the two FMAs,
$\alg{Fast2Sum}(b_{01}, b_{12})$ --- and \texttt{tw\_sum}, \texttt{tw\_prod},
\texttt{tw\_reci}, \texttt{tw\_div} likewise follow Algorithms 8, 9/10, 13
and~14. Algorithm~7 is not implemented at all, so the defect of
\S\ref{sec:roundtw} does not reach this code.

\paragraph{But the unguarded $\alg{Fast2Sum}$ of \S\ref{sec:delta2}(a) is
present verbatim,} in both Algorithm~18 and Algorithm~20:

\begin{quote}\ttfamily\small
two\_prod(a.x[0], i[0], \&z01\_plus, \&z01\_min);\\
fast\_two\_sum(a.x[1], z01\_plus, \&b0t, \&b1t);\quad /* $\alg{Fast2Sum}(b_1, z_{01}^{+})$ */\\
\dots\\
fast\_two\_sum(e1, e2, \&r.x[1], \&r.x[2]);\quad\ \ /* and the last line */
\end{quote}

\noindent
Those two routines are called by the reciprocal, the division \textbf{and} the
square root, so all three inherit the loss. Sorting the two arguments repairs
it at no arithmetic cost.

\paragraph{A second, related deviation.}
Algorithms 9 and 10 read
$(e_0,\dots,e_4) \leftarrow \alg{VecSum}(z_{00}^{+}, b_0, b_1, c, z_3)$, and
$\alg{VecSum}$ is built from $\alg{2Sum}$; the implementation uses
$\alg{Fast2Sum}$ at those steps instead. Same species of unproven shortcut.
(The \emph{addition} path does use $\alg{2Sum}$, correctly.)

\paragraph{And the source itself flags two errors in the appendix algorithms.}
In the routine computing $\frac32 - \Prodacc{2,3}(\bar b\,', i^{(1)})$ one
finds

\begin{quote}\ttfamily\small
// fault in algorithm specification, says z32 instead of z31 here\\
z3 = z31 + z01\_min;
\end{quote}

\noindent and, a few lines further down,

\begin{quote}\ttfamily\small
// Algorithms specifies * 0.5 somewhere, but For some reason the only way\\
// it is correct is by just doing it the normal way
\end{quote}

\noindent
These are the sign-folded variants (Algorithms 16, 17 and 19). We scoped them
out because they are presented as purely operational refinements --- they fold
a $+$ into a $-$ inside the multiplication, saving the explicit $2 - \cdot$ /
$\frac32 - \cdot$ and one operation, and are said to change no bound --- while
Theorems 9, 10 and 11 are stated about Algorithms 13, 14 and 15 as written in
the body. (We did formalise the two appendix algorithms the theorems genuinely
need, 18 and 20.)

That scoping now looks too generous. Algorithm~18 is an appendix algorithm too,
and it carried a real defect \emph{and} a $u^4$ constant more than twice the
announced one (\S\ref{sec:delta2}). The comments above say 16/17/19 as printed
do not work either. And these are the routines the implementation actually
runs --- so nothing in the verified chain covers the deployed reciprocal,
division or square root. \textbf{The appendix is worth a re-read, and is the
natural next thing to verify.}

\medskip
\noindent\emph{Scope: we read the main operation paths, not the constant-time
variants or the conversion routines.}

\section{In one line}

The published statements stand, with four amendments: \textbf{Algorithm 7's
inner test must be inverted}, \textbf{Algorithm 18's two $\alg{Fast2Sum}$ calls
must sort their arguments}, \textbf{$\delta_2$'s $u^4$ term is $620$, not
$260$} (hence Theorem~9's $1830$ / $1870$), and \textbf{Lemma 1 needs two
hypotheses}. Three further points (\S\ref{sec:thm6}, \S\ref{sec:thm8} and
Section 8.3's tail bound) are gaps in proofs whose statements are correct, and
Theorem~11 holds verbatim. The C reference implementation carries the
$\alg{Fast2Sum}$ defect unrepaired, and its own comments report two further
errors in the appendix algorithms (\S\ref{sec:impl}).

\end{document}
